
\subsection{Shared calibration and results}
The scan contains 1,499 slopes from $-2.499$ to $-1.001$
at spacing 0.001. Its best shared compromise is
\begin{equation}
\boxed{\alpha=-1.746,\qquad
\mathcal A=22.27648\,\mathrm{Gpc^{-3}\,yr^{-1}}.}
\end{equation}
The minimized joint score is 0.741507, below the acceptance value one.
The common median distance is $2.873356$ Gpc, respectively $21.71\%$
and $25.94\%$ below the public and high-cadence proxies. Rate discrepancies
remain as large as a factor $2.258$; passing the loose benchmark is not agreement
with its central values.
\begin{table}[H]\centering\small
\renewcommand{\arraystretch}{1.15}
\caption{Shared-LF rates. Raw rates include the optical-night convention; effective rates additionally include coverage 0.35. Rates are in $\mathrm{yr^{-1}}$.}
\begin{tabular}{@{}llrrrr@{}}
\toprule
Survey & Window & Raw & Effective & Rate/target & $q_{\rm med}$ \\\midrule
Public & Conservative & 6.576 & 2.302 & 1.151 & 1.4027 \\
High cadence & Conservative & 3.163 & 1.107 & 0.443 & 1.2832 \\
Public & Optimistic & 12.905 & 4.517 & 2.258 & 1.3144 \\
High cadence & Optimistic & 3.444 & 1.206 & 0.482 & 1.2586 \\
\bottomrule\end{tabular}\end{table}
\begin{table}[H]\centering\small
\renewcommand{\arraystretch}{1.15}
\caption{Viewing-angle diagnostics of the same shared model. Physical angles are in degrees; geometrical on-axis means $q<1$.}
\begin{tabular}{@{}llrrrr@{}}
\toprule
Survey & Window & $\theta_{\rm med}$ & $q<1$ & $q<1.5$ & $q>2$ \\\midrule
Public & Conservative & 8.037 & 25.71\% & 56.52\% & 24.14\% \\
High cadence & Conservative & 7.352 & 30.77\% & 64.38\% & 12.48\% \\
Public & Optimistic & 7.531 & 29.11\% & 63.78\% & 15.11\% \\
High cadence & Optimistic & 7.212 & 31.94\% & 66.10\% & 11.78\% \\
\bottomrule\end{tabular}\end{table}
\begin{table}[H]\centering\small
\renewcommand{\arraystretch}{1.15}
\caption{Detected luminosity medians, with the two luminosity definitions kept explicit.}
\begin{tabular}{@{}llrr@{}}
\toprule
Survey & Window & $L_\nu(t_{\rm dec})$ [erg s$^{-1}$ Hz$^{-1}$] & $\nu L_\nu(1\,\mathrm{d})$ [erg s$^{-1}$] \\\midrule
Public & Conservative & $1.13\times10^{34}$ & $2.95\times10^{45}$ \\
High cadence & Conservative & $7.57\times10^{32}$ & $1.97\times10^{44}$ \\
Public & Optimistic & $3.63\times10^{33}$ & $9.46\times10^{44}$ \\
High cadence & Optimistic & $6.49\times10^{32}$ & $1.69\times10^{44}$ \\
\bottomrule\end{tabular}\end{table}

The public-mode luminosity median is substantially higher than the high-cadence
one even though their distance medians are identical: the longer time between
required public visits selects intrinsically brighter afterglows. These
luminosity medians are predictions, not additional calibration targets.
All four detected means diverge in the ideal model at the selected slope;
the quoted medians remain finite.

\subsection{Accepted slopes, endpoint alternatives, and coverage}
At 0.001 resolution, shared calibrated slopes $-1.912\le\alpha\le-1.352$ pass both endpoints. Neighboring slopes $-1.913$ and $-1.351$ fail. This is a tolerance-defined feasible interval, not a confidence interval. The lower boundary follows the analytic distance constraint; the upper boundary comes from incompatible rate ratios. The $\alpha=-2$ model fails because its common median is only 2.275 Gpc, regardless of amplitude.
\begin{table}[H]\centering\small
\renewcommand{\arraystretch}{1.15}
\caption{Endpoint-specific fits use their own shared two-survey amplitude; they are separate alternatives, not a replacement of the four-case calibration. $\mathcal A$ has units $\mathrm{Gpc^{-3}\,yr^{-1}}$.}
\begin{tabular}{@{}lrrrr@{}}
\toprule
Calibration target & $\alpha$ & $\mathcal A$ & $D_{\rm med}$ [Gpc] & Target score \\\midrule
Both endpoints & -1.746 & 22.2765 & 2.8734 & 0.74151 \\
Conservative only & -1.613 & 18.8770 & 3.0784 & 0.59041 \\
Optimistic only & -1.727 & 20.1998 & 2.9060 & 0.71723 \\
\bottomrule\end{tabular}\end{table}
\begin{table}[H]\centering\small
\renewcommand{\arraystretch}{1.15}
\caption{Coverage sensitivity at the selected, fixed $\mathcal A$ and $\alpha$. C/O denote conservative/optimistic windows; rates are effective $\mathrm{yr^{-1}}$.}
\begin{tabular}{@{}rrrrrl@{}}
\toprule
Coverage & Public C & HC C & Public O & HC O & Passing endpoints \\\midrule
0.2 & 1.315 & 0.633 & 2.581 & 0.689 & Neither \\
0.35 & 2.302 & 1.107 & 4.517 & 1.206 & Both \\
0.5 & 3.288 & 1.581 & 6.452 & 1.722 & Conservative \\
\bottomrule\end{tabular}\end{table}

